\documentclass[10pt,a4paper]{amsart}
\usepackage[margin=19mm]{geometry}
\usepackage{amsmath,amssymb,mathtools}
\usepackage[scr=rsfs,cal=euler]{mathalfa}
\usepackage{xfrac}
\usepackage[unicode,hidelinks,bookmarks=false]{hyperref}
\newcommand{\ie}{i.e.\ }
\newcommand{\cf}{cf.\ }
\newcommand{\resp}{respectively\ }
\newcommand{\Subsection}{\subsection}
\providecommand{\nobreakdash}{\nobreak\hbox{-}}
\setlength{\emergencystretch}{2em}
\multlinegap0pt
\usepackage{amssymb}
\usepackage{mathtools}
\newtheorem{proposition}{Proposition}
\title{Effective dynamics and defect expansions for polynomial PDEs on thin annuli}
\author{Jean-Pierre Magnot}
\date{Source: arXiv:2602.12308v1 (CC BY 4.0)}
\begin{document}
\maketitle

\noindent\emph{Source and licence.} Effective dynamics and defect expansions for polynomial PDEs on thin annuli. Version arXiv:2602.12308v1 (CC BY 4.0), distributed by arXiv under the Creative Commons Attribution 4.0 International licence, http://creativecommons.org/licenses/by/4.0/, as recorded in the arXiv licence metadata for this exact version and shown on the abstract page https://arxiv.org/abs/2602.12308v1. Full licence notice: \texttt{licenses/arxiv-2602.12308-CC-BY-4.0.txt}.\par\smallskip

\noindent\textit{Editorial scope.} Counted unit: the article's proposition prop:compactness (source0.tex lines 697--710). The article's complete original proof of it is delivered with it (source0.tex lines 1427--1505, one \texttt{proof} environment of 79 lines, which the article places away from the statement). No mathematics is written, repaired or re-derived: the statement and the proof are the article's own lines, in the article's own notation, and no retained line is substituted or re-flowed.  Retained uncounted as the source-local context the proof needs: lines 1407--1410; lines 1414--1417.  Nothing was omitted from any retained span, so the package carries no omission record.  No retained line contains a citation, so the wrapper carries no bibliography and no bibliography entry.  No retained line contains a figure environment, an image-inclusion command or drawing code, so no image is delivered and the package declares no asset dependency.  Editorial formatting only, in addition: an AMS \texttt{amsart} preamble that restates the article's own declarations and macros used by the retained lines, namely the environments \texttt{proposition} on separate counters, together with the standard packages those lines need.  The retained environments print in this document as Proposition 1 in order of appearance, whereas the article prints the same items with the numbering of its own full document.  The complete native source of the article as distributed in the arXiv e-print for this exact version is bundled unchanged as \texttt{sources/194536/source0.tex}.
\begin{equation}
\label{eq:unif-bound-proof}
\sup_{t\in[0,T]}\|u_\varepsilon(t)\|_{\varepsilon,m}\le C(T),
\end{equation}

\begin{equation}
\label{eq:time-bound}
\|\partial_t u_\varepsilon\|_{L^2(0,T;H^{-m}(A_\varepsilon))}\le C(T),
\end{equation}

\begin{proposition}[Compactness]
\label{prop:compactness}
There exists a subsequence (still denoted $u_\varepsilon$) and a limit function
$u_0\in L^2(0,T;H^m(\mathbb S^1))$ such that
\[
u_\varepsilon \longrightarrow u_0
\quad\text{strongly in }L^2_{\mathrm{loc}}((-1,1)\times\mathbb S^1\times(0,T)),
\]
and
\[
\partial_\rho u_\varepsilon \longrightarrow 0
\quad\text{in }L^2(A_\varepsilon\times(0,T)).
\]
\end{proposition}

\begin{proof}[Proof of Proposition~\ref{prop:compactness}]
We prove radial rigidity and compactness.

\smallskip
\noindent\textbf{Step 1: Control of transverse derivatives.}
For \(m\ge1\), the definition of \(\|\cdot\|_{\varepsilon,m}\) includes the term
(with \(a=1,b=0\))
\[
\varepsilon^{2-(2m-1)}\int_{A_\varepsilon}|\partial_r u_\varepsilon|^2\,r\,dr\,d\theta
=
\varepsilon^{3-2m}\int_{A_\varepsilon}|\partial_r u_\varepsilon|^2\,r\,dr\,d\theta.
\]
In rescaled coordinates, \(\partial_r=\varepsilon^{-1}\partial_\rho\) and
\(r\,dr\,d\theta=\varepsilon(1+\varepsilon\rho)d\rho\,d\theta\), hence
\[
\int_{A_\varepsilon}|\partial_r u_\varepsilon|^2\,r\,dr\,d\theta
=
\varepsilon^{-1}\int_{-1}^1\!\!\int_{\mathbb S^1}
|\partial_\rho u_\varepsilon|^2(1+\varepsilon\rho)\,d\rho\,d\theta.
\]
Therefore, from~\eqref{eq:unif-bound-proof} we obtain
\[
\varepsilon^{3-2m}\cdot \varepsilon^{-1}
\int_{-1}^1\!\!\int_{\mathbb S^1}
|\partial_\rho u_\varepsilon|^2(1+\varepsilon\rho)\,d\rho\,d\theta
\le C(T)^2,
\]
i.e.
\begin{equation}
\label{eq:drho-bound}
\int_{-1}^1\!\!\int_{\mathbb S^1}
|\partial_\rho u_\varepsilon|^2\,d\rho\,d\theta
\le C(T)^2\,\varepsilon^{2m-2}.
\end{equation}
Since \(m\ge1\), the right-hand side tends to \(0\), which proves
\(\partial_\rho u_\varepsilon\to0\) in \(L^2((-1,1)\times\mathbb S^1\times(0,T))\).

\smallskip
\noindent\textbf{Step 2: Poincar\'e inequality in the transverse variable.}
For each fixed \((\theta,t)\), the 1D Poincar\'e inequality on \((-1,1)\) yields
\[
\int_{-1}^1\Bigl|u_\varepsilon(\rho,\theta,t)-\tfrac12\!\int_{-1}^1u_\varepsilon(\sigma,\theta,t)\,d\sigma\Bigr|^2\,d\rho
\le C\int_{-1}^1|\partial_\rho u_\varepsilon(\rho,\theta,t)|^2\,d\rho,
\]
with \(C\) independent of \(\varepsilon\).
Integrating in \(\theta\) and \(t\) and using~\eqref{eq:drho-bound} gives
\begin{equation}
\label{eq:transverse-avg}
\|u_\varepsilon-\bar u_\varepsilon\|_{L^2((-1,1)\times\mathbb S^1\times(0,T))}
\longrightarrow 0,
\end{equation}
where
\(\bar u_\varepsilon(\theta,t)=\tfrac12\int_{-1}^1 u_\varepsilon(\rho,\theta,t)\,d\rho\).

\smallskip
\noindent\textbf{Step 3: Compactness in \(\theta\).}
The uniform bound~\eqref{eq:unif-bound-proof} also yields control of tangential
derivatives up to order \(m\), hence
\[
\sup_{\varepsilon}\|\bar u_\varepsilon\|_{L^2(0,T;H^m(\mathbb S^1))}\le C(T).
\]
Together with the time regularity bound~\eqref{eq:time-bound} (or by Galerkin
finite-dimensionality), the Aubin--Lions lemma implies that \(\bar u_\varepsilon\)
is relatively compact in \(L^2(\mathbb S^1\times(0,T))\).
Thus, up to a subsequence,
\[
\bar u_\varepsilon \to u_0 \quad\text{strongly in }L^2(\mathbb S^1\times(0,T)).
\]

\smallskip
\noindent\textbf{Step 4: Strong convergence of \(u_\varepsilon\).}
Combining the strong convergence of \(\bar u_\varepsilon\) with
\eqref{eq:transverse-avg} yields
\[
u_\varepsilon \to u_0
\quad\text{strongly in }L^2_{\mathrm{loc}}((-1,1)\times\mathbb S^1\times(0,T)),
\]
and the limit is independent of \(\rho\).
\end{proof}

\end{document}
